\section{Related Work}\label{sec:related}
This section traces how pedestrian trajectory prediction has developed, identifies the limitations that remain in unstructured traffic, and states how each proposed component responds to them. The limitations are numbered L1--L4 and are summarised against the contributions in Table~\ref{tab:gaps}.

\subsection{From Social Models to Egocentric Prediction}
Early learning-based predictors treated a pedestrian as a sequence and modelled interaction between pedestrians. Social LSTM~\citep{alahi2016social} pooled the hidden states of neighbouring pedestrians, Social GAN~\citep{gupta2018socialgan} added adversarial training to produce several plausible futures, Social-STGCNN~\citep{mohamed2020stgcnn} replaced pooling by spatio-temporal graph convolutions, and Trajectron++-style models~\citep{salzmann2020trajectron} extended the generative formulation to heterogeneous agents. These models were developed for a fixed, bird's-eye view in which the positions of all agents are available, and their strength lies in modelling social interaction. They are not designed for an \emph{egocentric} setting, where a vehicle-mounted camera observes a pedestrian whose bounding box is influenced by the vehicle's own motion. Egocentric benchmarks, namely JAAD~\citep{rasouli2017jaad} and PIE~\citep{rasouli2019pie}, therefore added ego-vehicle speed and, later, optical flow as inputs, and PIE introduced a joint intention and trajectory baseline (PIETraj) built on recurrent networks. Surveys of the area~\citep{rasouli2020survey} note that the interaction between vehicle and pedestrian, and not only pedestrian--pedestrian interaction, is what matters for driving.

\textbf{Remaining limitation (L1).} Both families were developed and evaluated on data from structured traffic. When IDD-PeD was released, the error of established predictors on it was markedly higher than on PIE~\citep{iddped2025}, which indicates that cues that suffice under lane discipline, signals and crosswalks are insufficient when pedestrians walk along the carriageway, cross diagonally and hesitate.

\subsection{Transformer and Goal-Conditioned Predictors}
Attention-based models replaced recurrence to capture longer dependencies. Transformer forecasters for trajectories~\citep{giuliari2021transformer} and socio-temporal transformers such as AgentFormer~\citep{yuan2021agentformer} showed that self-attention is competitive with recurrent models. For egocentric prediction, MTN~\citep{yin2021mtn} used a multimodal transformer over bounding boxes, ego speed and optical flow with a deterministic decoder, and this combination is the strongest purely deterministic baseline reported on IDD-PeD~\citep{iddped2025}. A second line of work reduces uncertainty about the intended destination: endpoint-conditioned prediction~\citep{mangalam2020pecnet}, BiTraP~\citep{yao2021bitrap} with bidirectional goal estimation and SGNet~\citep{wang2022sgnet} with stepwise goal prediction. These stochastic models report the best of several sampled trajectories and obtain much lower errors on IDD-PeD than the deterministic models (Table~\ref{tab:prior})~\citep{iddped2025}, but best-of-$K$ numbers measure the quality of a set of hypotheses rather than of a single committed forecast, and the downstream planner still has to choose among them.

\textbf{Remaining limitation (L2).} Compact deterministic transformers rely on motion evidence alone. When a pedestrian pauses or changes direction, the motion history carries little information about the intended path, and no additional evidence about the \emph{situation} of the pedestrian is available to the model. Goal-conditioned methods address this by sampling several futures, not by supplying more evidence about the scene.

\subsection{Scene Context for Pedestrian Behaviour}
A natural response is to add scene context. Stacked recurrent models fused several contextual features for action anticipation, PedFormer~\citep{rasouli2022pedformer} modulates cross-modal attention and learns intention and trajectory jointly with gated multitask learning, and the dual-task model of Sharma et al.~\citep{sharma2025dual} fuses pedestrian regions of interest, scene attributes and past trajectories, refines cross-modal features with a progressive denoising attention module, trains with counterfactual manipulation of regions and segmentation maps, and predicts trajectories with a conditional variational autoencoder. These systems show that context helps for crossing intention and trajectory prediction. They also share three properties: context is injected by \emph{image-level features} (regions of interest, segmentation maps) that require an additional perception network; it is injected \emph{unconditionally}, so an uninformative or noisy context influences every sample; and the evaluation is on PIE and JAAD, i.e.\ structured traffic.

\textbf{Remaining limitation (L3).} It is not established how \emph{lightweight, native annotations} such as the intersection type, signalisation, road type, location type, motion direction and occlusion state provided by IDD-PeD can be injected into a compact attention model (i)~without altering the motion representation that is already learned and (ii)~without harming samples for which the context is uninformative. Context that cannot be switched off is a liability precisely in unstructured scenes, where the same label (for example an unsignalised junction) may or may not matter.

\subsection{Adaptive Fusion and Gating}
Selecting among computations per input has a long history in conditional computation. Hard and soft gates choose experts in sparsely-gated mixtures~\citep{shazeer2017moe}; discrete decisions are made trainable by the Gumbel-softmax relaxation~\citep{jang2017gumbel} and straight-through estimators~\citep{bengio2013ste}; and dynamic multimodal fusion~\citep{xue2022dynmm} chooses among fusion operations for each input of a multimodal model. These ideas have, to the authors' knowledge, not been applied to the fusion of static scene labels with a trajectory memory in egocentric pedestrian prediction, and their benefit when the amount of training data is small is untested; in the present study adaptive fusion adds parameters (Sec.~\ref{ssec:fusion}) without reducing computation, since all candidates are evaluated before one is selected.

\textbf{Remaining limitation (L4).} A fusion mechanism for context needs a guaranteed way to fall back to the motion-only prediction, and it needs to be evaluated under a matched protocol to establish whether the extra flexibility pays off.

\subsection{Parameter Sharing and Compact Models}
Models for pedestrian prediction are deployed on vehicles with limited memory and are trained on small datasets, so parameter efficiency matters. Sharing weights across layers reduces the parameter count and can regularise training: Universal Transformers~\citep{dehghani2019universal} repeat a single layer and ALBERT~\citep{lan2020albert} shares parameters across all layers. These studies share \emph{whole layers} (or all sub-modules) in large language-understanding models. In a multimodal encoder with two interacting streams, layers play different roles: the self-attention refines the same representation at every depth, whereas the source and cross attention exchange information between streams whose content changes with depth. It is therefore unclear which sub-modules can be tied safely, and sharing everything can harm accuracy (Sec.~\ref{sec:ablation}).

\textbf{Remaining limitation (L5).} A compact multimodal predictor has no principled way to reduce its parameters, and adding a context pathway increases them further.

\subsection{Positioning of the Present Contributions}
Table~\ref{tab:gaps} relates each limitation to the component that addresses it. The proposed architecture is deterministic, so that the effect of scene context and of parameter sharing can be isolated from the effect of a stochastic decoder; it is positioned as complementary to goal-conditioned sampling and not as a replacement. Native annotations replace image-level context features (L3), SAGE and the source mask of MASC provide an exact fallback to the motion-only prediction (L3, L4), ROUTE supplies a learned choice of fusion operation that is evaluated under a matched protocol (L4), MASC lets every decoder layer read the context through the attention it already uses (L2, L3), and SATIE offsets the added parameters (L5). Whether these remedies are sufficient is a matter of evidence, and Secs.~\ref{sec:results} and~\ref{sec:ablation} report where they did and did not help.

\begin{table*}[t]
\centering
\caption{Limitations of earlier approaches in unstructured traffic and the proposed components that address them.}
\label{tab:gaps}
\scriptsize
\begin{tabular}{p{0.3cm}p{4.9cm}p{4.4cm}p{5.5cm}}
\toprule
 & Limitation of earlier approaches & Where it arises & Proposed response\\
\midrule
L1 & Developed on structured-traffic data; performance drops on IDD-PeD & Social models, PIETraj, MTN~\citep{iddped2025} & Evaluation on IDD-PeD with native, unstructured-traffic annotations\\
L2 & Motion-only evidence is weak when pedestrians hesitate or turn & Deterministic transformers~\citep{yin2021mtn}; goal sampling gives hypotheses, not evidence & Native scene-context token added to the decoder memory (MASC)\\
L3 & Context needs image-level features and is injected unconditionally & PedFormer, dual-task CVAE model~\citep{rasouli2022pedformer,sharma2025dual} & Native labels; SAGE admission gate with exact fallback through the source mask (MASC)\\
L4 & Fusion operation is fixed; no fallback; benefit of adaptivity untested & Fixed fusion; dynamic fusion in other domains~\citep{xue2022dynmm} & ROUTE, adaptive fusion evaluated in a matched $2{\times}2$ study\\
L5 & Parameter growth; unclear which sub-modules can be tied & Whole-layer sharing~\citep{dehghani2019universal,lan2020albert} & SATIE: sharing of encoder self-attention only ($-6.07\%$)\\
\bottomrule
\end{tabular}
\end{table*}
